\chapter{System Architecture}
\label{ch:arch}

This chapter presents NatQuery's architecture as a diagram, then describes its design principles, its
components and the trust boundaries between them. Chapter~\ref{ch:how} follows one question through the
system in detail.

\section{Design Principles}

Four principles shape every part of the system (Table~\ref{tab:principles}). They follow from one
observation: a language model is a capable but untrusted author, so the safety of the system must not
depend on what the model writes.

\begin{table}[htbp]
\centering
\caption{Design principles and what they mean in the code.}
\label{tab:principles}
\small
\renewcommand{\arraystretch}{1.2}
\begin{tabularx}{\textwidth}{L{3.1cm} Y Y}
\toprule
\hdr \textbf{Principle} & \textbf{Meaning} & \textbf{Consequence} \\
\midrule
The model proposes, the system disposes & LLM output is text to be checked, never a command & The text the model writes is never executed; only SQL regenerated from the validated syntax tree is \\
Least visibility & A role sees only the tables and columns its policy grants & Hidden names are absent from the prompt and fail to resolve, exactly like names that do not exist \\
Defence in depth & No single control is trusted alone & Seven layers (Figure~\ref{fig:layers}); the database repeats the row scopes through row-level security \\
Everything is replaceable & Authentication, database, model, limiter and session store sit behind interfaces & A new login system, database engine or model is a new class, not a rewrite \\
Evidence by default & Every request leaves a record before it touches data & If the audit record cannot be written, the query does not run \\
\bottomrule
\end{tabularx}
\end{table}

\section{Architecture Diagram}

Figure~\ref{fig:arch} shows the path of one question. The numbered boxes are the stages a request passes
through in order; the dashed boxes are outside the trusted core (the user and the language model).
Stages 1--3 identify the user and screen the question, stages 4--7 build the prompt and obtain a
candidate query, stages 8--10 check and run it, and stages 11--13 clean up and answer. The orange arrow
is the repair loop: when the enforcer rejects a fixable query, the model receives a sanitised reason and
tries again, at most twice.

\begin{figure}[htbp]
\centering
\resizebox{\textwidth}{!}{% NatQuery architecture diagram (TikZ). Edit freely; compiled inside Chapter 2.
\begin{tikzpicture}[
  x=1cm, y=1cm,
  box/.style={draw=nqblue, fill=nqlight, rounded corners=2pt, align=center, minimum width=2.95cm, minimum height=1.15cm, font=\scriptsize, inner sep=2pt, text width=2.7cm},
  ext/.style={box, fill=white, dashed, draw=nqred},
  store/.style={cylinder, shape border rotate=90, aspect=0.22, draw=nqblue, fill=nqgrey, align=center, font=\scriptsize, minimum height=1.2cm, minimum width=2.4cm, inner sep=2pt},
  flow/.style={-{Latex[length=1.8mm]}, thick, nqblue},
  loop/.style={-{Latex[length=1.8mm]}, thick, nqorange},
  opt/.style={-{Latex[length=1.8mm]}, thick, dashed, gray},
  lab/.style={font=\tiny, inner sep=1pt, fill=white},
  tag/.style={circle, draw=nqgreen, fill=white, font=\tiny\bfseries, inner sep=0.8pt, text=nqgreen}
]
% ---- columns and rows
\def\cA{0}\def\cB{3.6}\def\cC{7.2}\def\cD{10.8}\def\cL{-4.7}
\def\rA{0}\def\rB{-2.5}\def\rC{-5.0}\def\rD{-7.5}\def\rE{-9.9}

% ---- row 1: identify and screen
\node[ext, minimum width=2.4cm, text width=2.1cm] (user) at (\cL,\rA) {User / CLI\\(asks in English)};
\node[box] (auth) at (\cA,\rA) {\textbf{1 Authentication}\\and session};
\node[box] (role) at (\cB,\rA) {\textbf{2 Role selection}\\highest priority or the one chosen};
\node[box] (guard) at (\cC,\rA) {\textbf{3 Input guard}\\and rate limit};
\node[store] (cache) at (\cD,\rA) {Schema cache\\+ semantics};
% ---- row 2: build the prompt
\node[ext, minimum width=2.4cm, text width=2.1cm] (llm) at (\cL,\rB) {LLM (external,\\untrusted)};
\node[box] (client) at (\cA,\rB) {\textbf{7 LLM client}\\SQL extraction, repair prompts};
\node[box] (prompt) at (\cB,\rB) {\textbf{6 Prompt builder}\\rules + schema + question};
\node[box] (select) at (\cC,\rB) {\textbf{5 Schema selection}\\rank, FK bridge, value hints};
\node[box] (view) at (\cD,\rB) {\textbf{4 Role-filtered}\\\textbf{schema view}};
% ---- row 3: enforce and execute
\node[box, fill=nqorange!25, draw=nqorange] (enf) at (\cA,\rC) {\textbf{8 SQL enforcer}\\parse, allow-list, resolve, row filters, masks, regenerate};
\node[box] (audit1) at (\cB,\rC) {\textbf{9 Audit record}\\``authorized'' (written before execution)};
\node[box] (exec) at (\cC,\rC) {\textbf{10 Executor}\\read-only transaction, \texttt{SET LOCAL ROLE}, session variables};
\node[store] (pg) at (\cD,\rC) {PostgreSQL\\+ row-level\\security};
% ---- row 4: return
\node[ext, minimum width=2.4cm, text width=2.1cm] (user2) at (\cL,\rD) {User / CLI\\(reads the answer)};
\node[box] (ans) at (\cA,\rD) {\textbf{13 Answer}\\rows, notes, reference id};
\node[box] (audit2) at (\cB,\rD) {\textbf{12 Audit record}\\``completed''};
\node[box] (rguard) at (\cC,\rD) {\textbf{11 Result guard}\\redaction, size caps};
\node[box, fill=white, dashed] (plan) at (\cD,\rD) {Plan and cost analysis\\(optional, \texttt{EXPLAIN})};
% ---- row 5: stores
\node[store] (pol) at (\cA,\rE) {Policy file\\(YAML)};
\node[store] (log) at (\cB+0.6,\rE) {Audit log\\(hash chain)};
\node[font=\scriptsize, align=left, anchor=west, text width=6.3cm] (leg) at (\cC-0.1,\rE) {%
  \tikz[baseline=-0.6ex]\draw[flow] (0,0)--(0.6,0); request path \quad
  \tikz[baseline=-0.6ex]\draw[loop] (0,0)--(0.6,0); repair loop\\[2pt]
  \tikz[baseline=-0.6ex]\node[tag]{P}; reads the policy \quad
  \tikz[baseline=-0.6ex]\node[tag]{A}; writes the audit log};

% ---- flow arrows
\draw[flow] (user) -- node[lab,above]{question} (auth);
\draw[flow] (auth) -- (role);
\draw[flow] (role) -- (guard);
\draw[opt] (cache) -- (view);
\draw[flow] (guard.south) -- ++(0,-0.55) -| node[lab, pos=0.25, above]{} (view.north);
\draw[flow] (view) -- (select);
\draw[flow] (select) -- (prompt);
\draw[flow] (prompt) -- (client);
\draw[<->, thick, nqred, dashed] (llm) -- node[lab,above]{prompt / SQL} (client);
\draw[flow] ([xshift=-0.35cm]client.south) -- node[lab, left]{SQL text} ([xshift=-0.35cm]enf.north);
\draw[loop] ([xshift=0.35cm]enf.north) -- node[lab, right]{reason} ([xshift=0.35cm]client.south);
\draw[flow] (enf) -- node[lab,above]{safe SQL} (audit1);
\draw[flow] (audit1) -- (exec);
\draw[flow] (exec) -- (pg);
\draw[<-, thick, nqblue] ([yshift=-0.25cm]exec.east) -- ([yshift=-0.25cm]pg.west);
\draw[flow] (exec) -- (rguard);
\draw[opt] (exec.south east) -- (plan.north west);
\draw[flow] (rguard) -- (audit2);
\draw[flow] (audit2) -- (ans);
\draw[flow] (ans) -- node[lab,above]{answer} (user2);

% ---- markers
\node[tag] at (role.north east) {P};
\node[tag] at (view.north east) {P};
\node[tag] at (enf.north east) {P};
\node[tag] at (exec.north east) {P};
\node[tag] at (auth.north east) {A};
\node[tag] at (audit1.north east) {A};
\node[tag] at (audit2.north east) {A};
\end{tikzpicture}
}
\caption{NatQuery architecture. Numbers give the order of the stages. P marks stages that read the
policy file; A marks stages that write an audit record.}
\label{fig:arch}
\end{figure}

\section{Components}

Table~\ref{tab:components} lists each component, where it lives in the code base and what it does. The
whole system is about 5,300 lines of Python in 57 modules.

\begin{table}[htbp]
\centering
\caption{Components of NatQuery and the packages that implement them.}
\label{tab:components}
\small
\renewcommand{\arraystretch}{1.15}
\begin{tabularx}{\textwidth}{L{2.6cm} L{4.3cm} Y}
\toprule
\hdr \textbf{Component} & \textbf{Package / file} & \textbf{Responsibility} \\
\midrule
Command line & \code{cli/} & The \code{natquery} command: connect, interactive shell, schema, policy, user and audit administration \\
Authentication & \code{auth/} & Provider interface; built-in local accounts (scrypt hashes, lockout); sessions with idle and absolute expiry \\
Input guard, limiter & \code{security/} & Heuristic screening of the question, per-user rate limits, secret redaction in logs \\
Schema store & \code{schema/} & Extracts the schema, caches it, merges descriptions, synonyms and value hints, ranks tables for large schemas \\
Policy, role view & \code{access/policy.py} \code{access/schema_view.py} & Loads the YAML policy, resolves role inheritance, builds the role-filtered schema \\
Prompt builder & \code{prompt/} & Assembles rules, role-visible schema and the question into model messages \\
LLM client & \code{llm/} & Model interface (Groq built in), SQL extraction and normalisation, repair prompts \\
SQL enforcer & \code{access/enforcer.py} \code{access/functions.py} & Parses the SQL, checks it, applies row filters and masks, regenerates the SQL \\
Repair loop & \code{orchestration/} & Orchestrates the stages; classifies errors so the repair prompt can give a targeted hint \\
Database backend & \code{backends/} & Connection pool, read-only execution, \code{SET LOCAL ROLE}, session variables, schema extraction (PostgreSQL) \\
Result guard & \code{access/result_guard.py} & Pattern redaction (for example card numbers) and size caps on returned values \\
Audit log & \code{audit/} & Append-only JSON lines, hash-chained, optional HMAC \\
Observability & \code{observability/} \code{execution/} & Cost and plan analysis with \code{EXPLAIN}, index suggestions for roles allowed to see plans \\
\bottomrule
\end{tabularx}
\end{table}

\section{Trust Boundaries and Defence in Depth}

Two things cross the boundary to the language model: the user's question and the role-visible schema
(table and column names, descriptions and value hints). Row data never does. Everything the model sends
back is untrusted. Figure~\ref{fig:layers} shows the seven layers between a question and the data, and
what each layer is there to stop.

\begin{figure}[htbp]
\centering
\resizebox{\textwidth}{!}{% Defence in depth: what each layer stops. Edit freely.
\begin{tikzpicture}[
  lay/.style={draw=nqblue, fill=nqlight, rounded corners=2pt, minimum width=5.2cm, minimum height=0.95cm, align=center, font=\small, text width=4.9cm},
  why/.style={font=\footnotesize, anchor=west, text width=8.3cm, align=left},
  flow/.style={-{Latex[length=2mm]}, thick, nqblue}
]
\def\dy{-1.25}
\node[draw=nqred, dashed, rounded corners=2pt, minimum width=5.2cm, minimum height=0.8cm, font=\small] (q) at (0,0) {Question (untrusted text)};
\foreach \i/\t/\w [count=\k] in {
  1/{Input guard, rate limit, session check}/{Obvious injection phrases, floods, expired or forged sessions},
  2/{Role-filtered schema view}/{The model never sees tables or columns the role may not use},
  3/{SQL enforcer}/{Anything the model still writes: hidden data, writes, dangerous functions, filter bypasses},
  4/{Read-only transaction, limits}/{Runaway queries, huge results, any write that slipped through},
  5/{Database role and row-level security}/{A bug in layers 2--3: the database limits the rows again},
  6/{Result guard}/{Secrets in values (card numbers, tokens) and oversized cells},
  7/{Audit log}/{Not a barrier but evidence: every decision is recorded and tamper-evident}} {
  \node[lay] (L\i) at (0,{\k*\dy}) {\textbf{\i}~~\t};
  \node[why] at ([xshift=0.4cm]L\i.east) {\w};
}
\node[draw=nqgreen, rounded corners=2pt, minimum width=5.2cm, minimum height=0.8cm, font=\small] (a) at (0,{8*\dy}) {Answer (only entitled rows)};
\draw[flow] (q) -- (L1);
\foreach \i [evaluate=\i as \j using int(\i+1)] in {1,...,6} {\draw[flow] (L\i) -- (L\j);}
\draw[flow] (L7) -- (a);
\end{tikzpicture}
}
\caption{Defence in depth: the seven layers between an untrusted question and the data.}
\label{fig:layers}
\end{figure}

\begin{table}[htbp]
\centering
\caption{What crosses each boundary.}
\label{tab:boundaries}
\small
\renewcommand{\arraystretch}{1.2}
\begin{tabularx}{\textwidth}{L{3.3cm} Y Y}
\toprule
\hdr \textbf{Boundary} & \textbf{Passes across} & \textbf{Never passes across} \\
\midrule
NatQuery $\rightarrow$ LLM provider & Question text; role-visible schema with descriptions and value hints & Row data, hidden tables or columns, identity attributes, passwords \\
LLM provider $\rightarrow$ NatQuery & One candidate SQL statement (untrusted) & Nothing is executed as received \\
NatQuery $\rightarrow$ database & The regenerated, validated SQL; the database role; identity as transaction settings & The model's original text \\
Database $\rightarrow$ user & Rows the role is entitled to, after redaction and caps & Secret columns (never granted), hidden rows \\
NatQuery $\rightarrow$ audit log & Who, which role, question, model SQL, executed SQL, tables, masks, repairs, outcome & Secrets (redacted before writing) \\
\bottomrule
\end{tabularx}
\end{table}

\FloatBarrier
